% Methods (Task 46, Scientific Reports): after the Discussion. v3 Threat Model and Methods plus Experimental Setup,
% with the Statistics subsection required by the journal and the LLM-use paragraph verbatim from v3 (last subsection).
% Every result number is a macro.
\section*{Methods}
\label{sec:methods}

\subsection*{Threat model}
\label{sec:threat}
The defender runs a flow classifier, an FT-Transformer \cite{gorishniy2021}, and a learned attack detector on top of
it (Fig.~\mainFigPipeline).
The attacker holds a malicious flow and wants the classifier to predict it as benign while the detector stays below its
threshold.
The attacker knows the weights of the classifier and of the detector (white box) and can change only the flow features
that a sender controls, within a per-feature budget and under the protocol relations of the constraint model below.
We consider three levels of attacker knowledge about the detector: none, a differentiable approximation of its scores
(stand-ins), and its exact score, with surrogate gradients only where the exact gradient is zero almost everywhere.
Following the adaptive-attack principle \cite{tramer2020,carlini2017}, a detector is judged by the attacker that models
it, and two detectors are compared at equal attacker effort.
Packet-level attacks, black-box attackers and attacks on the training data are out of scope.

\subsection*{Datasets and models}
\label{sec:setup}
We used LYCOS-IDS2017 \cite{rosay2021,rosay2022} and LycoS-Unicas-IDS2018 \cite{cantone2024} (called 2017 and 2018),
which re-extract the flows of CIC-IDS2017 and CSE-CIC-IDS2018 \cite{sharafaldin2018,aws_cse_cic_ids2018} after the
problems documented in the original datasets \cite{engelen2021,liu2022,lanvin2022}; both have the same \nFeatures{}
flow features (Supplementary Table~\suppTabData).
On 2017 we removed Heartbleed (\nHeartbleed{} rows) and the SQL injection web attack (\nSqlInjection{} rows, all after
the temporal cut), merged the remaining web attacks, removed feature vectors that carried more than one label and exact
duplicates, and split each (label, day) group in time with gaps between train, validation and test.
On 2018 we removed FTP-Patator (\nFtpRemoved{} rows) and DoS Slowhttptest (\nSlowhttpRemoved{} rows) because the two
classes were indistinguishable in the features, merged the DDoS and web attack variants, removed conflicting and
duplicate vectors, capped large classes, and split at random with stratification, because the file has no timestamps.
This left \nClassesSeventeen{} classes on 2017 and \nClassesEighteen{} on 2018; all cross-dataset comparisons use the
\nCommonClasses{} common attack classes.
We trained an FT-Transformer with \nSeeds{} seeds and an XGBoost reference with one recipe on both datasets, without
tuning on 2018 (Supplementary Note~\suppNoteImpl).
Test macro-F1 of the FT-Transformer was \fOneFTTSeventeen{} (seed range \fOneFTTSeventeenRange) on 2017 and
\fOneFTTEighteen{} (\fOneFTTEighteenRange) on 2018; XGBoost reached \fOneXGBSeventeen{} and \fOneXGBEighteen.
Clean explanations used up to \paramExplPerClass{} test rows per class that all seeds classified correctly, and attacks
and detectors used the first \paramPerClass{} of these rows per attack class (fewer when a class had fewer), the same
rows for every seed.

For every test row we computed the Euclidean distance to its nearest training row in the standardised feature space
and flagged rows with a distance below \paramNDThr.
We report three row sets: V0 (all source rows), V1 (flagged source rows removed in both datasets) and V2 (DoS Slowloris
and web attack also removed, \nVtwoClasses{} classes).
V1 is the primary set; V0 and V2 are sensitivity analyses.
The filter was applied after the attacks and after detector training, by removing source rows from the saved results,
and V1 was fixed as the primary set after the sensitivity results were seen, which is why V0 and V2 are reported as
well (Supplementary Note~\suppNoteNeardup).

\subsection*{Constraint model and validity checker}
\label{sec:constraints}
Of the \nFeatures{} features, \nControls{} are controls the attacker may change within a budget, \nDerived{} are
derived and recomputed from the controls, and \nImmutable{} are immutable (Supplementary Table~\suppTabConstraints),
following the problem-space view of constrained attacks \cite{pierazzi2020,alhussien2024}.
The budget of a control is $\epsilon$ times its clean value, or $\epsilon$ times the training median of its non-zero
values when the clean value is zero; class families whose mechanism is the timing and packet count (slow DoS, port
scan) have those controls frozen.
A separate validity checker, written apart from the derivation and projection that the attacks use, re-derived every
relation from the written specification and accepted every constrained adversarial row of every reported experiment
(for the adaptive attacks of the controls, \sbInvalid{} invalid of \sbRows{} rows; \axInvalid{} of \axRows{} for A8x;
\sxInvalid{} of \sxRowsWritten{} for A6x).
It shares the specification file and a few helpers with the attack code, so it shows internal consistency rather than
realisability.

\subsection*{Explanations and faithfulness}
For each FT-Transformer we computed three attention variants over the feature tokens.
The first is the row of the classification token in the last layer, averaged over heads; the second is the same row
averaged over layers and heads; the third is attention rollout \cite{abnar2020}.
We compared them with these references: IG with the training mean as baseline \cite{sundararajan2017},
Gradient\-SHAP \cite{lundberg2017}, occlusion of one feature at a time \cite{zeiler2014}, Tree\-SHAP of the XGBoost
model, the absolute standardised input $|x|$, and a random ranking.
Signed attributions were ranked by their absolute value.
Agreement between two explanations is the Spearman correlation of their feature rankings per row.
Faithfulness was measured by masking features in ranking order with the training mean and recording the probability
of the explained class, as the AUC of the deletion curve (lower is more faithful) and of the insertion curve (higher is
more faithful) \cite{petsiuk2018,samek2017}.
On adversarial rows we used the logit margin of the original class and the gap between most- and least-relevant-first
removal (AOPC gap, area over the perturbation curve) \cite{samek2017}, and subtracted the change measured under the
random control A0 at the same budget.
Removal-based scores can be confounded by the distribution shift of the masked inputs, can rank methods inconsistently
across evaluation strategies and can be statistically unreliable \cite{hooker2019,rong2022,tomsett2020}; our curves
mask with the training mean without retraining and share these caveats.
A parameter-randomisation check \cite{adebayo2018} is in Supplementary Note~\suppNoteRand.
All values are macro averages over classes and means over seeds.

\subsection*{Attacks}
\label{sec:attacks}
Attacks perturbed the controls of the raw flow within their budgets, recomputed derived features, and projected every
iterate back onto the constraint set; they used sign-gradient projected gradient descent with \paramSteps{} steps and
one random restart.
A0 draws one random point in the same box (random control), A1 maximises the benign logit margin (evasion), A1b targets
another attack class, A2 keeps the prediction and minimises the share of rollout attention on the clean
top-\paramTopK{} IG features, and A3 combines A1 with the A2 term ($\lambda \in \{100, 1000\}$ and a variant with
normalised gradients).
The budget $\epsilon$ ranged from $0.05$ to $2$; headline values use $\epsilon = \paramEpsHead$.
An unconstrained reference applied the same objectives in standardised space with an $L_\infty$ bound and no
constraints; its rows are not valid flows.

Every adaptive attacker maximises the benign margin, capped at $\kappa = \paramKappa$, minus a hinge on the detector
logit, in three settings ($\lambda \in \{1, 100\}$ and a normalised-gradient variant with weight \paramNormW), and
keeps the best projected iterate per row.
A6 attacks D, whose four scores are differentiable, so its objective is the exact score of D; it keeps the iterate
that is flipped and below the threshold under this score, and among those the one with the highest objective.
A8 attacks D$_+$ and replaces each consistency score by a differentiable stand-in: soft-rank Spearman correlation, a
soft top-\paramTopK{} Jaccard overlap, exact cosine similarity, and IG with \paramIGLoss{} steps inside the loss.
A8x attacks D$_+$ with its exact score wherever an exact gradient exists: attention, the four scores of D and the
cosine scores are computed exactly, and IG inside the loss uses \axIGSteps{} Gauss-Legendre steps, as in the true
scores.
Only the \nRankScores{} rank-based scores (Spearman and top-\paramTopK{} Jaccard) keep a surrogate gradient, soft ranks
whose temperature is annealed geometrically from \axTauStart{} to \axTauEnd{} over the steps, because their exact
gradient is zero almost everywhere.
At every step A8x evaluates the true D$_+$ score and keeps, per row, the flipped iterate with the lowest true score,
counting it as undetected only if it lies more than \axMargin{} below the threshold.
A6x is A6 with the selection rule of A8x: the same objective, settings, seeds and restart noise, with the iterate kept
by the true D score and the same margin.
It was added to test whether the different selection rules of A6 and A8x, rather than the consistency scores, explain
the comparison between D and D$_+$.
Detection was always evaluated with the true scores.
The labels follow the numbering of the code; A4, A5 and A7 are earlier variants that are not part of the analysis.

\subsection*{Detectors and controls}
\label{sec:detectors}
Detector D is a standardised, class-balanced logistic regression on four scores: maximum softmax probability, logit
margin, Mahalanobis distance and mean distance to the \paramKnn{} nearest benign training rows.
Detector D$_+$ adds \nConsistency{} consistency scores between attention (rollout, layer mean) and attribution (gradient
times input, IG with \paramIGDet{} steps), each measured by Spearman correlation, top-\paramTopK{} Jaccard overlap and
cosine similarity and oriented in the direction validated on validation data; its logistic regression is
L2-regularised.
Both detectors were trained only on validation data: \paramNegTrain{} benign validation rows predicted benign as
negatives and A1 rows generated on validation source rows as positives; the threshold was set at \paramFPR{} FPR on
\paramNegHeld{} held-out benign validation rows, and test negatives were \paramNegTest{} genuine benign test rows per
seed.

To test whether a difference between D and D$_+$ depends on attention, we trained \nControlsAll{} control detectors on
the same rows, with the same threshold rule and the inverse regularisation strength of D$_+$ (\apCplus).
Each keeps the four scores of D and replaces the \nConsistency{} consistency scores of D$_+$.
In a1, each row receives the attention of another row of the same true class (a random cyclic shift), so a1 is a
diagnostic control that no deployed system could compute; in a2, the attention is drawn from clean benign validation
rows that were used neither for training nor for the threshold; in b, attention is replaced by occlusion and
GradientSHAP; c0 is D with the regularisation of D$_+$ and no further score, c1 adds \nConsistency{} features of
standard normal noise, and c2 adds \nConsistency{} statistics of the logits and distances that use no explanation.
Each control was attacked with the same three settings by an attacker that models its own scores: exact gradients for
c0, c1 and c2, and the stand-ins of A8 for a1 and a2, with the selection rule of A6.
Control b was not attacked adaptively: its attack run exceeded the memory of the graphics card and was stopped before
any row was written.

\subsection*{Equal-effort evaluation, outcome rules and amendments}
\label{sec:protocol}
For each source row we counted an attack as successful and undetected if any candidate aimed at a detector flipped the
prediction to benign and stayed below that detector's threshold: D faced A1 and the three A6 settings, D$_+$ faced A1
and the three A8x settings, and each control faced A1 and the three settings of its own attacker.
The union with A1 and the three A8 settings for D$_+$ is reported for comparison, and the union with A6x for D tests
the selection rule.
This union models an attacker with query access to the detector; equal effort means the same candidates (restarts,
steps and settings) on both sides, not equal computation or equal optimiser quality.
The headline quantity is the difference D$_+$ minus D in the share of such source rows on V1 at
$\epsilon = \paramEpsHead$, per dataset; a positive difference means that D$_+$ left more rows undetected.
We also report the paired TPR difference at the \paramFPR{} FPR threshold on the flipped rows of each attack; for
attacks aimed at D$_+$ this difference is negative by construction and is not used as evidence on its own.

The outcome rules for the comparisons with A8x and with A6x were recorded before the runs, before the code of each
attacker was written, with the 2017 interval on V1 as the deciding quantity and 2018 reported without changing the
conclusion, because D alone already left few rows undetected there.
For A8x, an interval containing $0$ or lying above $0$ meant that the 2017 gain of D$_+$ measured with A8 did not hold
against A8x; for A6x, an interval above $0$ meant that the comparison does not depend on the selection rule.
Two amendments were recorded before any attack row of the attacker concerned existed: the tolerance of one sanity check
of A8x and the \axMargin{} selection margin, after single-precision noise failed the check at $10^{-5}$, and the chunk
size of A6x, set to that of A6 after a sanity check failed at a smaller chunk (Supplementary Table~\suppTabRuns).

\subsection*{Statistics}
\label{sec:statistics}
Each comparison is paired: the same source rows, seeds and negatives enter both detectors.
Per seed, V1 contains \nSrcVoneSeventeen{} source rows on 2017 (V0: \nSrcVzeroSeventeen) and \nSrcVoneEighteen{} on
2018, the same rows in every seed, and each detector was evaluated on \paramNegTest{} benign test negatives per seed.
For the equal-effort union the \paramCILevel{} interval of D$_+$ minus D is a paired percentile bootstrap over source
rows within seed (\paramBoot{} draws); for the paired TPR differences it is a percentile bootstrap pooled over seeds.
Neither interval captures variation between seeds, so the per-seed values are given next to every headline
difference, and means are reported with the range over seeds.
No hypothesis test was run, so no P value is reported; the comparisons are judged by their intervals and per-seed
signs, and the deciding rules were the outcome rules above.

\subsection*{Use of large language models}
\label{sec:llm}
During the preparation of this work the authors used Claude Opus 5.5 (Anthropic), through Claude Code, in order to write
and debug the analysis code and to draft and edit the manuscript text under the direction of the first author, and
Claude Sonnet 5.5 (Anthropic), used by the first author through the Claude app in a separate session, in order to review
intermediate results and to draft the instructions given to Claude Code.
After using these tools, the authors reviewed and edited the content as needed and take full responsibility for the
content of the published article.
The authors also checked the numerical results with scripts, and the first author checked them manually against the
result files.
The analysis code added in this version (the adaptive attacks on the control detectors, the attackers A8x and A6x
and their summaries) was also written by Claude Opus 5.5 through Claude Code under the direction of the first author.

Every number in the paper is generated from the result files by scripts and inserted as a macro, and the authors checked the numbers against the result files.
